% introduction.tex -- input by main.tex
\section*{Introduction}
Cell culture remains a central route by which viruses are detected and propagated, and the appearance of cytopathic effect (CPE) in an inoculated monolayer is treated in institutional guidance as a primary indicator of viral presence (ATCC Virology Culture Guide; CLSI M41-A, section 7.4.2; ASM CPE protocols). The same guidance instructs laboratories to replace growth medium containing approximately 10\% FBS with a maintenance medium containing approximately 2\% FBS at or near inoculation. The guidelines agree on the step but not on why it is taken: the reasons given range from serum interfering with viral attachment, to serum containing viral inhibitors, to slowing monolayer overgrowth and maintaining cell viability, and several guidelines give no reason at all. A step applied so consistently, with no agreed justification, raises the question of what the step itself does to the readout it precedes.

Under these recipes an Isolation culture differs from the pre-inoculation culture in two variables at once -- the inoculum and the serum concentration -- while a single outcome, CPE, is read. Serum deprivation is known to alter the morphology, growth and survival of cultured cells. If the step-down alone produces CPE-like morphology, CPE read under dual-medium conditions cannot by itself be attributed to the inoculum unless a matched uninoculated culture in the same maintenance medium is run alongside.


This study combines three lines of evidence. First, we describe how serum is set before and after inoculation in published Isolation methods and institutional recipes (R1). Second, we score whether those publications report Isolation-stage negative controls that match the maintenance medium (R2). Third, we test directly whether lowering FBS from 10\% to 2\% in uninoculated Vero E6 cultures yields morphology described as CPE by the contract laboratory (CRO) and by independent raters blind to culture condition (R3). Historical context on how dual media entered practice (P1) and literature context on the physical separation of virions from extracellular vesicles (P2) are summarised in the Discussion. \todoA{Decide whether P3 (genome-origin path census) enters this paper or a companion paper.}
